\documentclass[11pt]{article}
\usepackage[T1]{fontenc}
\usepackage{amsmath,amssymb}
\usepackage{cite}
\usepackage{booktabs}
\usepackage{graphicx}
\usepackage{xurl}
\usepackage{hyperref}
\hypersetup{
  colorlinks=true,
  linkcolor=blue,
  urlcolor=blue,
  citecolor=blue,
  pdfauthor={Jabulani Monchwe, Martin Chanza, Olorato Motitswane, Mfundo Monchwe}
}
\newcommand{\software}{\textsf{ReliableML}}
\newcommand{\code}[1]{\texttt{\small #1}}
\newcommand{\jabulaniauthor}{
\IEEEauthorblockN{Jabulani Monchwe}
\IEEEauthorblockA{\textit{Applied Business Analytics and Decision Making}\\
\textit{North-West University}\\
\textit{IDEA-Research Club Chairperson}\\
Johannesburg, South Africa\\
\href{mailto:monchwej.25@gmail.com}{monchwej.25@gmail.com}}
}
\newcommand{\martinchanzauthor}{
\IEEEauthorblockN{Martin Chanza}
\IEEEauthorblockA{\textit{Applied Business Analytics and Decision Making}\\
\textit{North-West University}\\
Mafikeng, South Africa\\
\href{mailto:martin.chanza@nwu.ac.za}{martin.chanza@nwu.ac.za}}
}
\newcommand{\oloratoauthor}{
\IEEEauthorblockN{Olorato Motitswane}
\IEEEauthorblockA{\textit{Eduvos Midrand Campus -- IT Faculty Data Science, South Africa}\\
\textit{IDEA-Research Club}\\
\href{mailto:oloratomotitswane@gmail.com}{oloratomotitswane@gmail.com}}
}
\newcommand{\mfundoauthor}{
\IEEEauthorblockN{Mfundo Monchwe}
\IEEEauthorblockA{\textit{Eduvos Midrand Campus -- IT Faculty Data Science, South Africa}\\
\textit{IDEA-Research Club}\\
\href{mailto:diditmfundo@gmail.com}{diditmfundo@gmail.com}}
}
\newcommand{\minenhleauthor}{
\IEEEauthorblockN{Minenhle Mchunu}
\IEEEauthorblockA{\textit{Applied Business Analytics and Decision Making}\\
\textit{Eduvos Midrand Campus -- IT Faculty Data Science, South Africa}\\
\textit{IDEA-Research Club}\\
\href{mailto:Mchunu.minenhle02@gmail.com}{Mchunu.minenhle02@gmail.com}}
}
\newcommand{\projecturl}{\url{https://github.com/Jabualani/ML-Reliabilty}}

\title{\software{} Method Note: Global and Local Explanations}
\author{\contactauthors}
\date{Version 0.1.1}
\begin{document}
\maketitle
\begin{abstract}
The tabular explainer compares a Monte Carlo permutation estimate of
interventional Shapley contributions with local weighted linear surrogate
weights in the style of LIME. Global summaries aggregate absolute
contributions over selected rows; local reports describe individual model
outputs.
\end{abstract}

\section{Inputs and outputs}
The fitted estimator must provide \code{predict}, \code{predict\_proba}, and
\code{classes\_}. Explanations require finite numeric features, at least two
columns, and reference/background data. Binary explanations default to the
positive class; multiclass use requires a selected class. \code{X} and
background data must use the same feature schema.

The global summary is mean absolute Shapley value by feature over explained
rows. Local results include signed Shapley contributions and LIME-style local
surrogate weights. Reason alignment compares these methods using rank
correlation, top-feature overlap, and sign agreement; bootstrap intervals can
summarize rank agreement. An optional user-supplied global importance mapping
is clearly distinguished from computed SHAP values. Optional exports include
input columns, predictions, labels when supplied, and per-row top
contributing features in CSV or supported serialized formats.

\section{Methods}
The Shapley approximation samples feature permutations and compares model
probabilities along each permutation against a reference background. Its
Monte Carlo error depends on sample count and feature dependence. The local
surrogate perturbs observations, weights nearby samples, and fits a local
linear explanation; fidelity depends on the perturbation and neighborhood
design. These are model-behavior explanations, not tests of causal effects.

\section{Interpretation and limits}
Agreement between attribution methods indicates explanation consistency
under their respective approximations. It does not demonstrate that a
prediction is correct, that features are valid reasons, or that no spurious
correlation is present. Explanations can be sensitive to background choice,
correlated features, perturbation realism, and model probability quality.
Use representative labeled holdout evaluation, subgroup/time analysis,
leakage checks, calibration, realistic counterfactuals, and domain review.
Causal claims require a defensible causal model and assumptions.

\section{Interface}
\begin{verbatim}
report = model.explain(X_test, class_label=1)
print(report.global_shap_importance)
print(report.local_explanations[0])
print(report.reason_alignment)
\end{verbatim}
An already-fitted external estimator can be passed to
\code{explain\_estimator}; labels are optional for explanations and required
only for the optional descriptive label evaluation.

\section{Further reading}
\begin{thebibliography}{3}
\bibitem{shap}
S. M. Lundberg and S.-I. Lee. A unified approach to interpreting model
predictions. \emph{Advances in Neural Information Processing Systems}, 2017.
\url{https://arxiv.org/abs/1705.07874}.
\bibitem{lime}
M. T. Ribeiro, S. Singh, and C. Guestrin. ``Why should I trust you?''
Explaining the predictions of any classifier. \emph{Proceedings of KDD},
2016. \url{https://doi.org/10.1145/2939672.2939778}.
\bibitem{reliability}
\software{} explanation implementation:
\url{https://github.com/Jabualani/ML-Reliabilty}.
\end{thebibliography}
\end{document}
